\paragraph{Effectiveness of Gated Delta Rule.}
In \cref{tab:strategy}, we compare several linear update strategies for the memory state. 
Directly replacing old memory with new information yields a Dice score of 74.68, indicating that a simple structural change in the update formula is insufficient. Introducing either $\boldsymbol{\alpha}_t$ or $\boldsymbol{\beta}_t$ alone slightly improves performance, suggesting that selectively forgetting old content or reinforcing new features helps capture dynamic cardiac boundaries. Combining $\boldsymbol{\alpha}_t$ and $\boldsymbol{\beta}_t$ further provides bidirectional control: $\boldsymbol{\alpha}_t$ governs the decay of stored shapes to remove noise or errors, and $\boldsymbol{\beta}_t$ sets the strength of newly written features to accommodate abrupt structural changes. 
As a result, GDR reaches a Dice score of 95.11, surpassing the baseline by about 1.7 percentage points.

In \cref{fig:wei:para_a,fig:wei:para_b}, $\boldsymbol{\alpha}_t$ and $\boldsymbol{\beta}_t$ concentrate at specific intervals. During stable frames, both decay and writing intensities remain moderate. In frames with blurred boundaries or heavy noise, $\boldsymbol{\alpha}_t$ increases forgetting to discard flawed information. In cases of sudden contour changes, $\boldsymbol{\beta}_t$ boosts new feature updates. Large gradient jumps, such as +40 or -30 in \cref{fig:wei:Grad}, illustrate the importance of robust gating. In \cref{fig:wei:pearson}, the correlation pattern forms block-like and diagonal structures, reflecting the coordinated adjustment of $\boldsymbol{\alpha}_t$ and $\boldsymbol{\beta}_t$. This mechanism refines echocardiographic boundaries, filters noise during complex cardiac motion, and retains useful features in stable periods for more precise segmentation.